% GENERATED FILE. Edit canonical lessons, then run npm run generate:books.
% canonical-source-hash: fnv1a64:af7ac56fd1c287e1
% canonical-lessons: KA-C90-grandson, KA-C90-granddaughter, KA-C90-relative, KA-C90-neighbours, KA-C90-uncle

\chapter{More of the Family}
\label{ch:ka-more-of-the-family}
\hlchaptermodality{90}

\begin{chapteropening}
\textbf{By the end of this chapter:} \emph{I can name a grandson, a granddaughter, a relative, the neighbours and an uncle.}

Complete the last lesson of chapter 90: I can name a grandson, a granddaughter, a relative, the neighbours and an uncle.
\end{chapteropening}

\section[\texorpdfstring{\kn{ಮೊಮ್ಮಗ} (mommaga)}{ಮೊಮ್ಮಗ (mommaga)}]{\kn{ಮೊಮ್ಮಗ} (mommaga) — grandson}

\label{lesson:KA-C90-grandson}



\begin{quote}
Before the new one: say the Kannada for let's go, then the Kannada for a holiday.
\end{quote}

\subsection*{You'll want to know: \kn{ಮೊಮ್ಮಗ}}
\textbf{\kn{ಮೊಮ್ಮಗ}} (\emph{mommaga}) — \textquotedblleft{}grandson\textquotedblright{}. It contains \textbf{\kn{ಮಗ}} (\emph{maga}), \textquotedblleft{}son\textquotedblright{}, which you already know: a grandson is a son's or daughter's son.

\begin{cousinweb}
\textbf{\kn{ಮೊಮ್ಮಗ}} (\emph{mommaga}) is native Dravidian, built on \emph{maga}, \textquotedblleft{}son\textquotedblright{}, and DEDR 4715 lists \emph{mom-maga} itself, \textquotedblleft{}grandson\textquotedblright{}; the entry also has Tamil \emph{marumakaṉ}.
\end{cousinweb}

\begin{grammarlens}[title={What you've built}]
A known word inside a new one.
\end{grammarlens}

\subsection*{Your turn}
\begin{itemize}
\raggedright
  \item \emph{Say it:} \emph{mommaga}
  \item \emph{Say it:} \emph{mommaga}, once more
  \item \emph{Say it:} say \emph{mommaga}
  \item \emph{Recall:} say the Kannada for let's go, then the Kannada for a holiday, then say \emph{mommaga} again
\end{itemize}

\begin{tcolorbox}[breakable,colback=teal!4,colframe=teal!35!black,title={Before you move on}]
What does \kn{ಮೊಮ್ಮಗ} mean? (\textquotedblleft{}Grandson\textquotedblright{}.) Say it once more.
\end{tcolorbox}
\section[\texorpdfstring{\kn{ಮೊಮ್ಮಗಳು} (mommagaḷu)}{ಮೊಮ್ಮಗಳು (mommagaḷu)}]{\kn{ಮೊಮ್ಮಗಳು} (mommagaḷu) — granddaughter}

\label{lesson:KA-C90-granddaughter}



\begin{quote}
Before the new one: say the Kannada for a holiday, then the Kannada for grandson.
\end{quote}

\subsection*{You'll want to know: \kn{ಮೊಮ್ಮಗಳು}}
\textbf{\kn{ಮೊಮ್ಮಗಳು}} (\emph{mommagaḷu}) — \textquotedblleft{}granddaughter\textquotedblright{}, built the same way on \textbf{\kn{ಮಗಳು}} (\emph{magaḷu}), \textquotedblleft{}daughter\textquotedblright{}.

\begin{cousinweb}
\textbf{\kn{ಮೊಮ್ಮಗಳು}} (\emph{mommagaḷu}) is native Dravidian, built on \emph{magaḷu}, \textquotedblleft{}daughter\textquotedblright{}, and DEDR 4715 lists \emph{mom-magaḷu} itself, \textquotedblleft{}granddaughter\textquotedblright{}; the entry also has Tamil \emph{marumakaṉ}.
\end{cousinweb}

\begin{grammarlens}[title={What you've built}]
The pair is complete.
\end{grammarlens}

\subsection*{Your turn}
\begin{itemize}
\raggedright
  \item \emph{Say it:} \emph{mommagaḷu}
  \item \emph{Say it:} \emph{mommagaḷu}, once more
  \item \emph{Say it:} say \emph{mommaga, mommagaḷu}
  \item \emph{Recall:} say the Kannada for a holiday, then the Kannada for grandson, then say \emph{mommagaḷu} again
\end{itemize}

\begin{tcolorbox}[breakable,colback=teal!4,colframe=teal!35!black,title={Before you move on}]
What does \kn{ಮೊಮ್ಮಗಳು} mean? (\textquotedblleft{}Granddaughter\textquotedblright{}.) Say it once more.
\end{tcolorbox}
\section[\texorpdfstring{\kn{ಸಂಬಂಧಿಕ} (sambandhika)}{ಸಂಬಂಧಿಕ (sambandhika)}]{\kn{ಸಂಬಂಧಿಕ} (sambandhika) — a relative}

\label{lesson:KA-C90-relative}



\begin{quote}
Before the new one: say the Kannada for grandson, then the Kannada for granddaughter.
\end{quote}

\subsection*{You'll want to know: \kn{ಸಂಬಂಧಿಕ}}
\textbf{\kn{ಸಂಬಂಧಿಕ}} (\emph{sambandhika}) — \textquotedblleft{}a relative\textquotedblright{}, from \textbf{\kn{ಸಂಬಂಧ}} (\emph{sambandha}), \textquotedblleft{}connection, relationship\textquotedblright{}.

\textbf{\kn{ಅವರು} \kn{ನಮ್ಮ} \kn{ಸಂಬಂಧಿಕರು}} (\emph{avaru namma sambandhikaru}) — \textquotedblleft{}they are our relatives\textquotedblright{}.

\begin{cousinweb}
\textbf{\kn{ಸಂಬಂಧಿಕ}} (\emph{sambandhika}) is built on Sanskrit \emph{sambandha}, \textquotedblleft{}connection, kinship\textquotedblright{} (CDIAL 12955), with the ending \emph{-ika}, \textquotedblleft{}one belonging to\textquotedblright{}: a kinsman.
\end{cousinweb}

\begin{grammarlens}[title={What you've built}]
The word for everyone the family tree connects.
\end{grammarlens}

\subsection*{Your turn}
\begin{itemize}
\raggedright
  \item \emph{Say it:} \emph{sambandhika}
  \item \emph{Say it:} \emph{sambandhika}, once more
  \item \emph{Say it:} say \emph{sambandhika}
  \item \emph{Recall:} say the Kannada for grandson, then the Kannada for granddaughter, then say \emph{sambandhika} again
\end{itemize}

\begin{tcolorbox}[breakable,colback=teal!4,colframe=teal!35!black,title={Before you move on}]
What does \kn{ಸಂಬಂಧಿಕ} mean? (\textquotedblleft{}A relative\textquotedblright{}.) Say it once more.
\end{tcolorbox}
\section[\texorpdfstring{\kn{ನೆರೆಯವರು} (nereyavaru)}{ನೆರೆಯವರು (nereyavaru)}]{\kn{ನೆರೆಯವರು} (nereyavaru) — the neighbours}

\label{lesson:KA-C90-neighbours}



\begin{quote}
Before the new one: say the Kannada for granddaughter, then the Kannada for a relative.
\end{quote}

\subsection*{You'll want to know: \kn{ನೆರೆಯವರು}}
\textbf{\kn{ನೆರೆಯವರು}} (\emph{nereyavaru}) — \textquotedblleft{}the neighbours, the people next door\textquotedblright{}, from \textbf{\kn{ನೆರೆ}} (\emph{nere}), \textquotedblleft{}nearby, next door\textquotedblright{}.

\begin{cousinweb}
\textbf{\kn{ನೆರೆಯವರು}} (\emph{nereyavaru}) is native Dravidian: \emph{nere}, \textquotedblleft{}proximity, neighbourhood\textquotedblright{}, which DEDR 3770 lists itself, and \emph{avaru}, \textquotedblleft{}those people\textquotedblright{} (DEDR 1). The neighbours are \textquotedblleft{}the people close by\textquotedblright{}.
\end{cousinweb}

\begin{grammarlens}[title={What you've built}]
The people who are close without being family.
\end{grammarlens}

\subsection*{Your turn}
\begin{itemize}
\raggedright
  \item \emph{Say it:} \emph{nereyavaru}
  \item \emph{Say it:} \emph{nereyavaru}, once more
  \item \emph{Say it:} say \emph{nereyavaru}
  \item \emph{Recall:} say the Kannada for granddaughter, then the Kannada for a relative, then say \emph{nereyavaru} again
\end{itemize}

\begin{tcolorbox}[breakable,colback=teal!4,colframe=teal!35!black,title={Before you move on}]
What does \kn{ನೆರೆಯವರು} mean? (\textquotedblleft{}The neighbours\textquotedblright{}.) Say it once more.
\end{tcolorbox}
\section[\texorpdfstring{\kn{ಮಾವ} (māva)}{ಮಾವ (māva)}]{\kn{ಮಾವ} (māva) — uncle (mother's brother); father-in-law}

\label{lesson:KA-C90-uncle}



\begin{quote}
Before the new one: say the Kannada for a relative, then the Kannada for the neighbours.
\end{quote}

\subsection*{You'll want to know: \kn{ಮಾವ}}
\textbf{\kn{ಮಾವ}} (\emph{māva}) — \textquotedblleft{}uncle\textquotedblright{}, in particular the mother's brother. The same word also means \textquotedblleft{}father-in-law\textquotedblright{}: in the south the two roles were often the same person.

\begin{cousinweb}
\textbf{\kn{ಮಾವ}} (\emph{māva}) is native Dravidian, one word with Tamil \emph{māmaṉ} and Telugu \emph{māma} (DEDR 4813). DEDR lists \emph{māva} itself, \textquotedblleft{}mother's brother; spouse's father\textquotedblright{}. DEDR 4813 compares Sanskrit \emph{māma(ka)}.
\end{cousinweb}

\begin{grammarlens}[title={What you've built}]
One word, two relations.
\end{grammarlens}

\subsection*{Your turn}
\begin{itemize}
\raggedright
  \item \emph{Say it:} \emph{māva}
  \item \emph{Say it:} \emph{māva}, once more
  \item \emph{Say it:} say \emph{māva}
  \item \emph{Recall:} say the Kannada for a relative, then the Kannada for the neighbours, then say \emph{māva} again
\end{itemize}

\begin{tcolorbox}[breakable,colback=teal!4,colframe=teal!35!black,title={Before you move on}]
What does \kn{ಮಾವ} mean? (\textquotedblleft{}Uncle (mother's brother); father-in-law\textquotedblright{}.) Say it once more.
\end{tcolorbox}
